% 04_analysis_kor_use_vs_share.do가 생성. 직접 고치지 말 것.
활용도 1과 지표 A의 Spearman 순위상관 & 0.80 \\
활용도 2와 지표 B의 Spearman 순위상관 & 0.86 \\
(참고) 활용도 1과 활용도 2 & 0.97 \\
(참고) 지표 A와 지표 B & 0.96 \\
상위 10 겹침: 활용도 1 -- A & 6 \\
상위 30 겹침: 활용도 1 -- A & 19 \\
상위 10 겹침: 활용도 2 -- B & 5 \\
상위 30 겹침: 활용도 2 -- B & 19 \\
\(|G_1|\) 평균 & 18.0 \\
\(|G_1|\) 최댓값 & 74 \\
\(|G_1| \ge 20\)인 소분류 수 & 48 \\
\(|G_2|\) 평균 & 15.0 \\
\(|G_2|\) 최댓값 & 60 \\
\(|G_2| \ge 20\)인 소분류 수 & 39 \\
\(G_1\)과 취업자 수의 Spearman 순위상관 & 0.90 \\
\(G_2\)와 총 주업 근로시간의 Spearman 순위상관 & 0.87 \\
\(G_1\)과 \(G_2\)의 Spearman 순위상관 & 0.96 \\
\(\mathrm{var}(\ln \mathrm{A})\) & 4.05 \\
\quad 활용도 1 분산의 몫 & 1.52 \\
\quad 취업자 수 분산의 몫 & 0.52 \\
\quad 공분산(\(2\,\mathrm{cov}\))의 몫 & -1.04 \\
\(\ln\) 활용도 1과 \(\ln E\)의 Pearson 상관 & -0.59 \\
\(\ln \mathrm{A}\)를 \(\ln E\)에 회귀한 기울기(규모 탄력성) & 0.00 \\
\quad 같은 회귀의 \(R^2\) & 0.00 \\
\(\mathrm{var}(\ln \mathrm{B})\) & 5.97 \\
\quad 활용도 2 분산의 몫 & 1.45 \\
\quad 총 주업 근로시간 분산의 몫 & 0.35 \\
\quad 공분산(\(2\,\mathrm{cov}\))의 몫 & -0.80 \\
\(\ln\) 활용도 2와 \(\ln H\)의 Pearson 상관 & -0.56 \\
\(\ln \mathrm{B}\)를 \(\ln H\)에 회귀한 기울기(규모 탄력성) & -0.14 \\
\quad 같은 회귀의 \(R^2\) & 0.01 \\
표본 30 미만 소분류 수: 활용도 1 상위 10 중 & 2 \\
표본 30 미만 소분류 수: 지표 A 상위 10 중 & 0 \\
표본 30 미만 소분류 수: 활용도 2 상위 10 중 & 3 \\
표본 30 미만 소분류 수: 지표 B 상위 10 중 & 0 \\
표본 30 미만 소분류 수(129개 중) & 6 \\
취업자 수 최댓값 / 최솟값(배) & 1,282 \\
활용도 1과 취업자 수의 Spearman 순위상관 & -0.56 \\
지표 A와 취업자 수의 Spearman 순위상관 & -0.01 \\
활용도 2와 취업자 수의 Spearman 순위상관 & -0.54 \\
지표 B와 취업자 수의 Spearman 순위상관 & -0.10 \\
